% ESEMPIO 9: callout con codice e linguetta, callout con più puntatori, personaggio con bolla.
%   Passo 1: callout con del codice, che punta a un nodo, con una linguetta col titolo (stile box label).
%   Passo 2: un solo callout con più code verso più nodi (stile callouts, definito in examples/fig_callouts.tex;
%            serve la libreria TikZ decorations.pathreplacing, caricata in main.tex).
%   Passo 3: personaggio (tikzpeople) con bolla di dialogo.
%
%   Codice colorato con minted (come nel deck ICWS di Nicola): non è abilitato qui perché serve Pygments
%   installato (nel PC e nella CI) e shell-escape (già attivo in .latexmkrc). Per usarlo:
%     preambolo di main.tex:  \usepackage{minted}   \setminted{breaklines, breakanywhere}
%     nel callout, al posto del testo a larghezza fissa:  \parbox{.45\textwidth}{\inputminted{json}{file.json}}

\begin{frame}
    \frametitle{Example: callouts, code, several pointers, characters}

    \begin{adjustbox}{width=.95\linewidth, max height=.8\textheight, center}
        \begin{tikzpicture}
    % Stili per un callout con PIÙ puntatori (da tex.stackexchange.com, domanda 432276): store path salva il
    % percorso del nodo, callouts={(punto1), (punto2), ...} disegna una "coda" verso ciascun punto.
    \tikzset{
        % senza coda: rettangolo senza riempimento, serve a disegnare il callout con più code (vedi sotto)
        basic question no callout/.style={
            rectangle,
            preaction={draw=none, thick},
            fill=none
        },
        add path/.style={
            decoration={show path construction,
                moveto code={\xdef\savedpath{\savedpath (\tikzinputsegmentfirst)}},
                lineto code={\xdef\savedpath{\savedpath -- (\tikzinputsegmentlast)}},
                curveto code={\xdef\savedpath{\savedpath .. controls (\tikzinputsegmentsupporta) and (\tikzinputsegmentsupportb) ..(\tikzinputsegmentlast)}},
                closepath code={\xdef\savedpath{\savedpath -- cycle}}
            },
            decorate
        },
        store path/.style={add path},
        store path/.prefix code={\xdef\savedpath{}},
        callouts/.style={
            store path,
            append after command={
                foreach \target in {#1}{
                    ($(callout)!2pt!-90:\target$)--\target --($(callout)!2pt!90:\target$)
                } \savedpath
            },
            alias=callout
        }
    }

    \node[rboxes] (a) at (0,0) {Element A};
    \node[rboxes] (b) at (4,0) {Element B};
    \node[rboxes] (c) at (8,0) {Element C};

    % Passo 1: callout con del codice dentro e una linguetta col titolo (stile box label). Un solo puntatore:
    % callout absolute pointer. Il codice è testo a larghezza fissa; con minted (vedi sotto) avrebbe i colori.
    \node[basic question no color, callout absolute pointer={(a.north)}, anchor=south,
        visible on=<1>] (code callout) at ([yshift=1.3cm]a.north east) {%
        \parbox{.32\textwidth}{\ttfamily\scriptsize
            \{\\
            ~~"key": "value",\\
            ~~"status": 1\\
            \}}%
    };
    \node[box label, visible on=<1>] at (code callout.north) {\textbf{Label tab}};

    % Passo 2: un unico callout con tre code, verso tre nodi (stile callouts, definito sopra).
    % Il disegno (bordo e riempimento) è applicato solo al passo 2 con on slide: così anche le code, che fanno
    % parte dello stesso percorso, compaiono solo a quel passo.
    \path[basic question no callout, on slide={<2> preaction={draw, very thick}, fill=VeryLightGrey!50}]
        node[callouts={(a.south), (b.south), (c.south)}, visible on=<2>] at (4,-2.2) {%
            \parbox{.3\textwidth}{One bubble with several pointers}%
        };

    % Passo 3: personaggio (tikzpeople) con bolla di dialogo; l'ancora "mouth" è la bocca del personaggio.
    \node[devil, minimum size=.6cm, visible on=<3>] (attacker) at (10.5,-1.8) {};
    \node[speaking question, callout absolute pointer={(attacker.mouth)}, anchor=east,
        visible on=<3>] at ([xshift=-.6cm, yshift=-.3cm]attacker.west) {%
        \parbox{.2\textwidth}{A speech bubble from a character}%
    };

\end{tikzpicture}

    \end{adjustbox}

\end{frame}
